\documentclass[11pt,a4paper]{article}
\usepackage[utf8]{inputenc}
\usepackage[T1]{fontenc}
\usepackage[spanish,es-nodecimaldot]{babel}
\usepackage{lmodern}
\usepackage{geometry,booktabs,tabularx,array,enumitem,amsmath,graphicx,float,caption,hyperref}
\geometry{top=2.3cm,bottom=2.3cm,left=2.4cm,right=2.4cm}
\hypersetup{colorlinks=true,linkcolor=black,urlcolor=blue}
\setlength{\parindent}{0pt}
\setlength{\parskip}{0.55em}
\setlength{\tabcolsep}{4pt}
\renewcommand{\arraystretch}{1.15}
\setlist{leftmargin=*,topsep=3pt,itemsep=3pt}
\newcolumntype{Y}{>{\raggedright\arraybackslash}X}
\emergencystretch=2em
\raggedbottom

\title{\textbf{Exploración de modelos de inteligencia artificial multimodal\\
para la identificación de macromicetos mediante dispositivos móviles}\\[0.6em]
\normalsize Resultados de las semanas 9 y 10: integración de clima y sustrato}
\author{Daniel Esteban Camacho Ospina\\Cesar Augusto Pulido Cuervo}
\date{}

\begin{document}
\maketitle

\begin{abstract}
Se evalúa el aporte incremental de geolocalización, clima y sustrato a la
clasificación de diez géneros de macromicetos. El ensayo conserva las
predicciones de un clasificador visual MobileNetV3-Large previamente entrenado
y ajusta tres clasificadores de contexto mediante regresión logística. La
fusión se realiza sobre las probabilidades, con un peso fijo de 0,1 para el
contexto. En las 3\,452 observaciones de validación de FungiTastic
\texttt{v002}, el control visual obtuvo un F1 macro de 0,971674; la variante
con clima alcanzó 0,971950 y la variante con sustrato, 0,971880. Los intervalos
de las diferencias incrementales incluyen cero y no permiten establecer una
mejora concluyente. La evaluación externa de 4\,895 observaciones muestra que
transferir directamente el contexto geográfico puede reducir la exactitud.
Como el conjunto externo carece de clima y sustrato compatibles, esa evaluación
no mide su aporte externo con datos completos. El trabajo corresponde a un
ensayo de fusión tardía con la rama visual congelada, no a un ajuste conjunto
de la red multimodal ni a una evaluación final independiente.
\end{abstract}

\section{Objetivo y alcance}
Las semanas 9 y 10 se dedicaron a evaluar si el clima y el sustrato aportan
información adicional a la imagen y la ubicación. Se conserva la secuencia
incremental del proyecto: A, imagen; B, imagen y geolocalización; C, las
modalidades anteriores más clima; D, las anteriores más sustrato. El modelo
basado únicamente en coordenadas, implementado como diagnóstico en un commit
anterior, no se identifica aquí como B: omitir la imagen impediría atribuir la
diferencia B--C exclusivamente a la incorporación del clima.

El ensayo utiliza los recursos disponibles en este equipo: metadatos de
entrenamiento y validación, el paquete climático oficial y predicciones
visuales previamente guardadas. Se entrenaron nuevas ramas tabulares en CPU,
manteniendo congelado el modelo visual. Esta estrategia permite estudiar una
fusión controlada, pero no sustituye el entrenamiento conjunto de una red
multimodal en el equipo con GPU.

El informe mantiene el nivel taxonómico del experimento: la salida es uno de
diez \emph{géneros}, no una identificación definitiva de especie. Hábitat,
elevación, fecha, identificadores y etiquetas taxonómicas no son entradas de
los clasificadores de contexto de este ensayo.

\section{Datos y protocolo experimental}
\subsection{Cohortes y referencia visual}
Se utilizaron 17\,799 observaciones de entrenamiento y 3\,452 de validación
de FungiTastic \texttt{v002}, conservando las particiones oficiales. La
referencia visual es MobileNetV3-Large S5, configuración C, seleccionada
previamente por F1 macro en validación. La denominación \emph{S5-C} identifica
ese checkpoint previo y no debe confundirse con el modelo C de modalidades.
Es el mismo checkpoint visual en la validación interna y en el conjunto externo.

Las probabilidades visuales ya estaban disponibles y no se recalcularon. Se
comprobaron los identificadores, las etiquetas, el orden de clases y las
cohortes antes de reutilizarlas. Para cada observación se promedian las
probabilidades de sus fotografías, de modo que las observaciones con varias
vistas cuentan una sola vez en el ajuste tabular y la evaluación.

El F1 macro por observación es la métrica principal: promedia el F1 de los diez
géneros con el mismo peso. Se reportan también exactitud, aciertos y errores.
La validación ya había intervenido en la selección del modelo visual; por
tanto, estos resultados no son independientes de todo el proceso de desarrollo.
No se ejecutó una nueva evaluación en el test interno.

\subsection{Codificación y ajuste del contexto}
Las coordenadas válidas con calidad \texttt{OK} se redondean a 0,1 grados y se
representan mediante tres componentes sobre la esfera, con una marca de
ausencia. La resolución angular no equivale a una distancia fija ni constituye
por sí sola una garantía de privacidad. El clima se representa con
\texttt{Bio1}--\texttt{Bio19}, estandarizadas con estadísticas de entrenamiento
y acompañadas de indicadores de ausencia. El sustrato se codifica mediante
categorías aprendidas solo en entrenamiento, con una categoría desconocida.
La información de calidad se utiliza para decidir qué valores son admisibles;
no se incorpora como una categoría predictiva independiente.

Las tres ramas de contexto se ajustan por separado mediante regresión
logística multiclase, con ponderación equilibrada de clases, regularización
\(C=1\), optimizador L-BFGS, tolerancia \(10^{-5}\) y un máximo de 1\,000
iteraciones. El escalado y el vocabulario de categorías se obtienen solo de
las observaciones de entrenamiento. El procedimiento no realiza búsqueda de
hiperparámetros en validación ni en el conjunto externo.

\subsection{Fusión y estimación de incertidumbre}
Sean \(p_A(y\mid x)\) las probabilidades visuales agregadas de una observación
y \(q_M(y\mid z_M)\) las de la rama de contexto del modelo
\(M\in\{B,C,D\}\). La salida combinada es
\begin{equation}
 p_M(y\mid x,z_M)=(1-\alpha)p_A(y\mid x)+\alpha q_M(y\mid z_M),
 \qquad \alpha=0,1.
\end{equation}
El peso se fijó antes de calcular los resultados y no se modificó para
favorecer alguna variante. Cuando no existe contexto utilizable se toma
\(\alpha=0\), recuperando exactamente la predicción visual. La combinación
convexa conserva la normalización de las probabilidades; no implica que
estén calibradas.

Se calcularon diferencias B--A, C--B y D--C, con intervalos del 95\,\% mediante
2\,000 remuestreos pareados de observaciones y semilla 42. Los pares conservan
la misma observación en ambos modelos. Los intervalos describen la variación
muestral de esta comparación; no incorporan variación entre semillas de
entrenamiento ni corrigen el uso previo de validación.

\section{Semana 9: incorporación del clima}
\subsection{Fuente, disponibilidad y limitación temporal}
Se descargó y verificó el paquete oficial \texttt{climaticData.zip}, de
545\,926\,654 bytes. El descargador original solicita \texttt{climatic.zip},
pero el servidor publica el primer nombre. Se comprobaron la integridad CRC y
el SHA-256 del paquete, y se conservaron las fuentes y sus hashes.

El producto empleado es \textbf{CHELSA}, no ERA5-Land. Se utilizan las 19
variables bioclimáticas promedio de 1981--2010 descritas por FungiTastic
\cite{fungitastic}. Este cambio respecto a la propuesta se hace explícito:
el ensayo mide contexto climático de largo plazo, no las condiciones
meteorológicas de la fecha de captura ni el microclima del espécimen.
Los valores se conservan en su escala almacenada; no se presentan como grados
Celsius o milímetros sin una conversión verificada. No se incluye humedad
relativa. Las series mensuales descargadas se preservan, pero no se usan aquí.

\begin{table}[H]
\centering
\caption{Disponibilidad del contexto en la cohorte común de FungiTastic.}
\label{tab:coverage}
\begin{tabular}{lrrrr}
\toprule
Partición & Observaciones & GPS \texttt{OK} & Clima completo & Sin clima \\
\midrule
Train & 17\,799 & 17\,795 & 17\,782 & 17 \\
Validación & 3\,452 & 3\,452 & 3\,442 & 10 \\
\bottomrule
\end{tabular}
\end{table}

La unión se realiza por identificador de observación y conserva los registros
sin clima, con ausencia explícita, para mantener cohortes idénticas. Hay 745
observaciones de entrenamiento anteriores a 2011; las normales de 1981--2010
pueden contener información posterior a esos eventos. En consecuencia, el
análisis es retrospectivo y no reproduce estrictamente la información
climática que habría estado disponible en cada fecha.

\subsection{Resultados incrementales}
@@VALIDATION_TABLE@@

A y B presentan la misma exactitud y F1 macro en validación. C corrige dos
errores de B e introduce uno, para un saldo de un acierto adicional. El aumento
de F1 macro es pequeño y su intervalo incluye cero
(tabla~\ref{tab:deltas}). No se establece una mejora concluyente al añadir el
clima con esta codificación, arquitectura tabular y peso de fusión.

@@DELTA_TABLE@@

\section{Semana 10: incorporación del sustrato}
\subsection{Auditoría y armonización}
El sustrato describe el soporte sobre el que crece el hongo; no se equipara
con el hábitat. En los metadatos originales, 1\,811 observaciones de train y
una de validación carecen de \texttt{substrate} detallado y recurren a
\texttt{metaSubstrate=wood}. La regla de curaduría las marca
\texttt{LOW\_PRECISION}. Esa marca indica menor detalle, no un error de
etiqueta demostrado. Las tasas son 10,1747\,\% y 0,0290\,\%, respectivamente.

Para que el modelo no distinga madera detallada y madera sin especificación
como si fueran soportes biológicos diferentes, se agrupan madera viva, madera
muerta y madera no especificada en una categoría amplia común. Las categorías
restantes se conservan según la curaduría y el vocabulario se ajusta solo en
train. Los valores desconocidos o inadmisibles se representan como tales.
No se eliminaron observaciones ni se modificaron etiquetas después de ver las
predicciones.

El desequilibrio de detalle proviene de los CSV de origen y de una partición
temporal. El desglose por año y género confirma su presencia en varios años,
incluidos registros recientes de train; no permite atribuirlo únicamente a la
antigüedad ni a un cambio de formulario no documentado. Con un solo caso de
baja precisión en validación no es posible estimar de forma fiable el desempeño
de ese estrato por separado.

\subsection{Resultados y sensibilidad}
D obtiene el mismo número de aciertos que C. Corrige dos errores de C e
introduce dos distintos; su F1 macro disminuye ligeramente, de 0,971950 a
0,971880. El intervalo de D--C incluye cero, por lo que no hay evidencia de
una mejora consistente por incorporar sustrato en este ensayo.

@@SENSITIVITY_TABLE@@

Las perturbaciones se aplican por observación, sin separar sus vistas. La
ausencia total recupera el control visual por construcción. En algunos
escenarios, retirar GPS o clima produce un F1 mayor que el de las entradas
completas; por tanto, los resultados no respaldan una mejora monótona al
acumular modalidades. La similitud del F1 al permutar sustrato no demuestra
que las probabilidades sean idénticas ni que esa variable resulte inútil en
otras formas de fusión. Estas son pruebas de sensibilidad y no simulaciones
validadas de disponibilidad real en Android.

\section{Evaluación externa y cambio de dominio}
\subsection{Conjunto y condiciones}
Se reutilizaron las predicciones externas del mismo checkpoint visual S5-C y
se verificó nuevamente el SHA-256 de las 9\,649 fotografías del conjunto
\texttt{ext\_full\_001}. Se evaluaron 4\,895 observaciones: 3\,348 de
Colombia y 1\,547 de MIND.Funga. Dentro del grupo colombiano, 309 corresponden
a la Orinoquía. Los grupos se contienen entre sí y no deben sumarse ni tratarse
como pruebas independientes.

Las fuentes externas no se usaron para ajustar coeficientes, el peso de fusión
ni hiperparámetros. Sin embargo, ya habían sido evaluadas en una fase anterior
del proyecto. Se trata de una auditoría de robustez sobre fuentes conocidas,
no de un test externo nuevo e intacto.

Los metadatos externos no contienen las variables CHELSA ni sustrato compatible
con este ensayo. Se dejaron ausentes, sin sustituirlos por datos inferidos de
la etiqueta. De las 4\,895 observaciones, 2\,964 tienen GPS de calidad
\texttt{OK}; MIND.Funga no aporta ubicación. Esta evaluación mide transferencia
geográfica y comportamiento con modalidades faltantes: \textbf{no prueba el
aporte externo de clima y sustrato completos}.

\subsection{Transferencia directa y filtro geográfico}
@@EXTERNAL_TABLE@@

En Colombia, aplicar directamente B o C reduce la exactitud del control
visual de 88,32\,\% a 87,54\,\%. D alcanza 88,23\,\%, todavía por debajo de A.
En la Orinoquía, A obtiene 77,02\,\%, B/C 76,05\,\% y D 77,02\,\%.
La transferencia de asociaciones geográficas aprendidas principalmente en
Dinamarca no puede suponerse beneficiosa en Colombia. La diferencia D--C
externa no se interpreta como beneficio del sustrato: ese campo está ausente y
las ramas tabulares se ajustaron con distintos conjuntos de entradas.

También se evaluó un filtro previo de dominio geográfico, definido mediante
los percentiles 0,5 y 99,5 de latitud y longitud de train. Fuera de esos límites
se desactiva la contribución contextual. Ninguna observación externa cumple
simultáneamente los límites; por ello las variantes filtradas coinciden con A.
Esta igualdad es una propiedad de la regla y no una mejora predictiva ni una
validación de un detector de dominio. El filtro no se ajustó con etiquetas
externas y no garantiza cobertura adecuada de otras regiones.

\section{Disponibilidad para uso móvil}
La disponibilidad en el dataset no implica disponibilidad automática en una
aplicación. Foto y fecha pueden obtenerse localmente; GPS requiere permiso y
una localización utilizable. El clima derivado de una fuente remota se
incorpora al sincronizar, salvo que exista un paquete local compatible. No se
presupone que la persona conozca sustrato o hábitat: deben ser campos opcionales,
con un estado desconocido explícito.

Las rutas de ausencia evaluadas preservan la inferencia visual cuando no hay
contexto utilizable. Aún no se ha comprobado la equivalencia de la inferencia
en Android, su latencia, memoria, consumo ni la adquisición efectiva de los
campos en un teléfono. Tampoco se deben intercambiar directamente las
normales CHELSA con variables ERA5-Land sin adaptar y validar el modelo.

\section{Conclusiones y trabajo siguiente}
El ensayo de fusión tardía muestra variaciones pequeñas en validación: el
clima produce un acierto neto adicional y el sustrato no aumenta la exactitud
respecto a la variante climática. Los intervalos incrementales de F1 incluyen
cero. Estos resultados no justifican afirmar una mejora estable ni descartar
de forma general el potencial de esas modalidades.

La auditoría externa evidencia que el contexto geográfico puede perjudicar
la transferencia y que es necesario definir el comportamiento ante fuentes
incompletas o regiones distintas. La falta de clima y sustrato externos impide
una conclusión sobre su aporte multimodal completo fuera de FungiTastic.

El trabajo siguiente comprende:
\begin{enumerate}
\item Repetir A/B/C/D mediante ajuste conjunto en GPU y varias semillas,
manteniendo las mismas cohortes y controles. Comparar esos resultados por
separado con el ensayo de rama visual congelada.
\item Obtener contexto climático externo compatible y sustrato con
procedencia verificable antes de atribuir una mejora multimodal en Colombia.
\item Diseñar la evaluación prospectiva con información climática disponible
en la fecha de cada evento y con un conjunto independiente previamente definido.
\item Revisar errores por género, clases pequeñas y vistas discordantes, sin
corregir etiquetas o excluir ejemplos únicamente por su predicción.
\item Verificar adquisición de metadatos, rutas de ausencia y costo de
inferencia en un dispositivo Android de referencia.
\end{enumerate}

\section*{Trazabilidad y reproducción}
La configuración fija, los modelos tabulares, las métricas por condición y
los intervalos se conservan bajo \path{reports/s9_s10/frozen_context/}.
Las auditorías de procedencia, calidad y revisión visual permanecen en
\path{reports/s9/}. Los parámetros de cada rama se guardan en JSON junto con
su codificador y estadísticas de escalado. Las 21 pruebas automatizadas
aprobadas comprueban, entre otros aspectos, ajuste exclusivo en train,
modalidades admitidas, ausencias, alineación de cohortes y reanudación del
entrenamiento con datos sintéticos; no sustituyen una validación predictiva.

El informe se genera desde las salidas registradas, sin volver a entrenar:
\begin{quote}\small\ttfamily
python scripts/reports/build\_s9\_s10\_report.py
\end{quote}
Se entrega una fuente LaTeX autónoma y el PDF compilado. Las tablas se
construyen directamente desde los CSV y JSON de resultados.

\begin{thebibliography}{9}
\bibitem{fungitastic}
L. Picek y colaboradores.
\emph{FungiTastic: A multi-modal dataset and benchmark for image categorization}.
arXiv:2408.13632, sección 2.1.
\url{https://arxiv.org/html/2408.13632v2}.
\end{thebibliography}
\end{document}
